% IWAI 2026 submission --- draft outline
% Based on the Springer LLNCS samplepaper.tex template.
%
\documentclass[runningheads]{llncs}

\usepackage[T1]{fontenc}
\usepackage{graphicx}
\usepackage{booktabs}
\usepackage{amsmath}
\usepackage{amssymb}
\usepackage{hyperref}
\usepackage{color}
\renewcommand\UrlFont{\color{blue}\rmfamily}
\urlstyle{rm}

\begin{document}

\title{An Active Inference Agent for Atari Pong\\
       with Closed-Form Variational Inference}
\titlerunning{Active Inference on Atari Pong}

\author{Pablo Bustos\inst{1}\orcidID{0000-0000-0000-0000}}
\authorrunning{P. Bustos}
\institute{Robolab, Universidad de Extremadura, C\'aceres, Spain\\
\email{pbustos@unex.es}}
\maketitle

\begin{abstract}
We build the smallest closed-form Active Inference agent that we
could make work on a real Arcade Learning Environment game ---
\texttt{ALE/Pong-v5}, RAM observations --- and use the build as a
vehicle for two contributions. First, an end-to-end specification of
a gradient-free AIF agent: pretrained Variational Bayesian Gaussian
Sums (VBGS) for ball and opponent dynamics via closed-form CAVI on
30\,000 frames of free play, a $K{=}2$ mixture Kalman belief filter
with bounce-aware physics, EFE planning by deterministic horizon-3
tree search, and a streaming recurrent mixture model (rMM) for
online aim policy. Second, an honest empirical record: nineteen
paired-seed interventions, of which seven were positive and twelve
were not, evaluated at $N{=}100$ episodes per arm with paired
$t$-tests. The agent reaches a mean score of $+1.94$ from a
rule-based baseline of $-14$. The substantive lesson is not the
absolute score: it is that the two largest improvements
(+4.8 and +5.3 points per episode) came not from the generative
model but from clean velocity tracking through bounces and from
replacing a sampled MPPI planner by deterministic tree search. We
identify a consistent pattern across negative results --- smoothing
and stale-prediction interventions lose, responsiveness
interventions win --- that we conjecture generalises beyond Pong.
We position our numbers against AXIOM Gameworld~10k and Atari~100k
RL baselines (BBF, DreamerV3, EfficientZero) as reference points,
not direct competitors, and discuss the apples-to-oranges issues
explicitly.

\keywords{Active Inference \and Expected Free Energy \and Variational
Bayes \and Mixture Models \and Atari \and Pong \and CAVI \and
Negative Results.}
\end{abstract}

% =========================================================================
\section{Introduction}
% =========================================================================

\paragraph{Motivation.} Reinforcement learning has dominated Atari
benchmarks for over a decade, but at the cost of large sample
complexity, opaque value functions, and limited interpretability.
Active Inference (AIF), grounded in the Free Energy Principle, offers
an alternative in which behaviour emerges from minimising Expected
Free Energy (EFE) under a generative model of the world. Recent
results such as AXIOM~\cite{ref_axiom} demonstrate that AIF agents
with object-centric VBGS components, online structure growth, and
Bayesian model reduction can be competitive with sample-efficient RL
on a custom 10-game benchmark. What is missing from the literature
is a stripped-down counterpart: \emph{the simplest possible
closed-form AIF agent that still works on a real ALE game}, together
with an honest account of which interventions during its development
helped and which did not.

\paragraph{What this paper is, and what it is not.}
This paper is \emph{not} a benchmark contribution. We do not claim
to beat AXIOM, BBF, DreamerV3, or EfficientZero on any standard
metric, and we are explicit about the input-modality and pretraining
caveats that would make any such claim ill-founded
(Section~\ref{sec:comparison}). What this paper offers instead is:
\begin{enumerate}
  \item A \emph{minimal} fully closed-form AIF agent for
        \texttt{ALE/Pong-v5} --- pretrained VBGS dynamics on RAM
        observations, conjugate NIW priors, CAVI updates,
        Kalman-style belief filtering, EFE planning by tree
        search, online rMM aim policy. No gradients are used
        anywhere in the pipeline.
  \item An honest empirical record built around a paired-seed
        bench harness ($N{=}100$ episodes per arm, paired
        $t$-tests): seven positive interventions, twelve negative
        ones, all documented at the same level of detail.
  \item Two non-modelling lessons that we believe generalise:
        (i) clean velocity tracking through bounces and walls
        dominates any architectural change to the generative
        model itself; (ii) a deterministic horizon-3 tree search
        beats MPPI ($H{=}15$, $N{=}128$) by $+5.3$ points on the
        same code path. Both findings are isolated by the paired
        bench, and neither is a result that AXIOM, BBF or
        DreamerV3 could have surfaced.
\end{enumerate}

\paragraph{Outline.} Section~\ref{sec:related} surveys AIF agents and
sample-efficient RL on Atari. Section~\ref{sec:background} fixes
notation for FEP, EFE, and the conjugate Bayesian updates used.
Section~\ref{sec:model} specifies the generative model.
Section~\ref{sec:agent} describes the belief filter, action selector,
and the recurrent mixture model for online aim policy.
Section~\ref{sec:methodology} states the experimental protocol.
Section~\ref{sec:results} traces the improvement ladder.
Section~\ref{sec:negatives} catalogues the negative results.
Section~\ref{sec:comparison} compares to AXIOM and RL baselines.
Section~\ref{sec:discussion} concludes.

% =========================================================================
\section{Related Work}\label{sec:related}
% =========================================================================

\subsection{Active Inference Agents.}
[TODO --- cite Friston FEP review; Da Costa AIF tutorial; Tschantz
et al.\ scaling AIF; Markovic et al.\ on planning under EFE; the
AXIOM paper~\cite{ref_axiom}; VBGS structure learning. Position
this work as the simplest possible closed-form AIF agent that still
beats a hand-coded baseline on a standard ALE game.]

\subsection{Sample-efficient RL on Atari.}
The Atari 100k benchmark~\cite{ref_atari100k} (100\,000 environment
steps, $\approx 2$ hours of human gameplay) is the closest RL
counterpart to the sample budget regime we operate in. Leading
methods include EfficientZero~\cite{ref_efficientzero} (MuZero with
self-supervised auxiliaries; 194\% mean human-normalised score),
BBF~\cite{ref_bbf} (model-free SR-SPR with scaled networks), and
DreamerV3~\cite{ref_dreamerv3} (world-model agent). On
\texttt{ALE/Pong-v5} specifically, DQN-class methods reach $+20$ at
roughly 10M frames.

\subsection{Object-centric and structured world models.}
AXIOM~\cite{ref_axiom} is the most direct comparator: object-centric
generative model with VBGS components, expanded online via Bayesian
model reduction, evaluated on a 10-game Gameworld~10k benchmark in
which the game \emph{Bounce} is the Pong analog. Our agent shares the
VBGS choice and closed-form CAVI updates but is single-object and
uses pretraining rather than online structure growth.

% =========================================================================
\section{Background}\label{sec:background}
% =========================================================================

\paragraph{Free Energy and Expected Free Energy.}
[TODO --- 1 page. Variational free energy $F = \mathrm{KL}[q\|p] -
\log p(o)$; expected free energy decomposition into pragmatic value
$\mathbb{E}_q[\log p(o^*)]$ and epistemic value
$\mathbb{H}[p(s|o)] - \mathbb{E}_q\mathbb{H}[p(s|o)]$.]

\paragraph{Conjugate updates and CAVI.}
[TODO --- NIW prior on Gaussian components, posterior in closed form,
mixture responsibilities by digamma terms, CAVI as alternating E/M.
Cross-reference PONG.md \S2--\S4.]

% =========================================================================
\section{Generative Model}\label{sec:model}
% =========================================================================

\paragraph{State and observation.}
6-D latent state $s = [b_x, b_y, v_x, v_y, p_y, o_y]$ (ball position,
ball velocity, player paddle, opponent paddle); 4-D observation
$o = [b_x, b_y, p_y, o_y] / 255$ read from RAM bytes
$[49, 54, 51, 50]$. Velocities are latent and recovered by the filter.
Action set $\{0, 2, 3\}$ = \{NOOP, UP, DOWN\}.

\paragraph{Ball dynamics.}
A VBGS with $K=4$ components fitted by CAVI on 30k frames of
free play, with a Normal--Inverse--Wishart prior. Wall bounces are
captured by separate mixture components; the conditional predictor
\texttt{predict\_conditional}$(s_t \to s_{t+1})$ uses the
component-mixture posterior.

\paragraph{Opponent dynamics.}
Symmetric VBGS over $(\text{ball state}, o_y, o_{y, t+1})$;
empirically the opponent in ALE/Pong follows a fixed lag controller
$\approx 12$\,px above the ball.

\paragraph{Likelihood.}
Linear-Gaussian, $p(o\,|\,s) = \mathcal{N}(Hs, R)$, with $H$ selecting
$(b_x, b_y, p_y, o_y)$ and $R = 10^{-4} I$. $R$ is optionally
re-estimated online from Kalman innovations.

\paragraph{Prior preferences $C$.}
Three Gaussian-mixture goals encode (i) ball on the left half, (ii)
player paddle intercepts the ball, (iii) opponent paddle far from the
ball. The intercept term shrinks $\sigma_{p_y}$ adaptively as the
ball approaches the player's $x$.

% =========================================================================
\section{The Agent}\label{sec:agent}
% =========================================================================

\subsection{Belief Filter}
$K{=}2$ mixture Kalman with the VBGS as the conditional transition;
process noise $Q = \mathrm{diag}(10^{-4},10^{-4},10^{-4},10^{-4},10^{-3},10^{-4})$.
An EMA frame-difference velocity estimator (parameter $\alpha=0.7$)
runs in parallel with the Kalman estimate to bypass the post-contact
velocity spike (Kalman $v_x$ jumps to $\pm 0.3$ after paddle contact;
the EMA does not).

\subsection{Bounce-aware velocity}
The single largest engineering win in this paper (Section~\ref{sec:results}).
Two heuristics fire inside the EMA update:
\begin{itemize}
  \item Sign-flip bounce: $|\Delta v| > 0.003$ and
        $\mathrm{sgn}(\Delta v) \neq \mathrm{sgn}(v_{\text{EMA}})$
        $\Rightarrow$ hard-reset (no blending).
  \item Wall-clip bounce: if $b_y$ is within $0.005$ of the
        top/bottom edge for one frame, $\Delta v_y \approx 0$ is
        suppressed; treat as a bounce.
\end{itemize}

\subsection{EFE Action Selection}
Deterministic horizon-3 tree search over the 3-action space (27 leaves).
Depth-0 uses the full VBGS \texttt{predict} (bounce-aware via the
$K{=}2$ mixture); depths 1--2 use a linear-Gaussian rollout matrix
$F$ for speed. The target observation $o^*$ is computed once per
frame from the current belief and shared across all action branches
(per-branch recomputation introduces VBGS noise that degrades
selection).

\subsection{Recurrent Mixture Model (rMM) for online aim policy}
A streaming AXIOM-style rMM over continuous features
$(t_y, o_y^{\text{arrival}})$ with NIW + Beta-head components:
deterministic CRP growth ($K_{\max}=16$, log-likelihood spawn
threshold $-2.0$). The planner sweeps 9 continuous aim candidates
within the reachable window and selects by Bayes-UCB on the
Beta-head win-rate. Legacy 2D Beta-grid checkpoints auto-migrate
into the rMM on load.

% =========================================================================
\section{Experimental Methodology}\label{sec:methodology}
% =========================================================================

\paragraph{Paired-seed bench.}
Two arms run with identical seeds, identical pretrained VBGS, and
identical warm rMM state. Each episode uses
\texttt{seed} $= \texttt{SEED\_BASE} + \texttt{ep}$. Each arm is
evaluated for $N=100$ episodes; per-seed differences feed a paired
$t$-test. The protocol is implemented in
\texttt{bench\_lock\_slack.py}.

\paragraph{Metrics.}
Per-episode score (sum of $\pm 1$ rewards over 21 points), contacts
per episode, no-contact losses (rallies in which the paddle never
touched the ball), wins ($\text{score} > 0$), and contact-rally win
rate.

\paragraph{Significance.}
We treat $|\Delta| < 1$ point as inside the noise band at $N=100$
(standard error on the paired mean is $\approx 0.7$--$1.0$). At
$N=20$ paired arms, the noise band widens to $\approx 2$ points and
we required $p < 0.001$ to claim a positive result.

% =========================================================================
\section{Results: An Improvement Ladder}\label{sec:results}
% =========================================================================

Table~\ref{tab:ladder} traces the chronology of substantive changes.
The full set of nineteen attempted interventions is in
Section~\ref{sec:negatives}.

\begin{table}[t]
\caption{Improvement ladder. Mean score per episode, paired-seed
bench, $N{=}100$.}\label{tab:ladder}
\centering
\begin{tabular}{lrr}
\toprule
Configuration & Mean score & $\Delta$ \\
\midrule
Rule-based baseline                          & $-14.4$ &       \\
Tracking current $b_y$ ($H{=}3$)             & $-20.9$ & $-6.5$ \\
$+\,$ landing-prediction target              & $-10.8$ & $+10.1$ \\
$+\,$ EMA raw velocity                       & $-8.4$  & $+2.4$  \\
$+\,$ bounce-detection + wall-clip           & $-1.7$  & $+6.7$  \\
$+\,$ tree-search planner (replaces MPPI)    & $-1.7$  & $+5.3$\textsuperscript{*} \\
$+\,$ rMM re-featured on $t_y$               & $+1.9$  & $+3.6$  \\
$+\,$ \texttt{late\_bounce\_strategic} ($\alpha{=}0.6$) & $+2.6$  & $+0.7$  \\
\bottomrule
\end{tabular}
\smallskip\\
\textsuperscript{*}MPPI$(H{=}15)$ vs tree$(H{=}3)$ on the same code path.
\end{table}

\paragraph{Two non-modeling wins.}
The largest single improvement was bounce-aware velocity tracking
($+4.8$ mean per episode, paired $t=4.99$, $p<0.001$). The second
was diagnosing that the production code used MPPI$(H{=}15, N{=}128)$
while the bench harness used tree search $(H{=}3)$. MPPI's averaging
washed out the per-input bounce-detection signal; long-horizon
linear-Gaussian rollouts drifted past depth $\sim 5$. Together these
two changes account for $>90\%$ of the total agent improvement,
\emph{neither of which is a change to the generative model itself}.

\paragraph{Strategic targeting.}
\texttt{late\_bounce\_strategic} (\S\ref{sec:agent}) was the first
intervention that exploited a behavioural model of the opponent ---
the observation that the hand-coded ALE/Pong opponent tracker lags
on direction reversals near its own paddle. An $\alpha$-sweep at
$N=100$ identified $\alpha{=}0.6$ as the optimum
($\Delta = +0.71$, $p \approx 0.05$).

% =========================================================================
\section{Negative Results}\label{sec:negatives}
% =========================================================================

Twelve out of nineteen documented interventions had no effect or
actively regressed the score. Reporting them is part of this paper's claim
that paired-seed evaluation is necessary for closed-form AIF
research, where every change to the generative model or the planner
plausibly looks right in isolation.

\begin{table}[t]
\caption{Negative results, all $N{=}100$ paired-seed bench. ``EMA''
denotes a smoothing intervention; ``anticip.''\ denotes an
intervention that commits to a stale prediction.}\label{tab:neg}
\centering
\begin{tabular}{lrl}
\toprule
Intervention & $\Delta$\,mean & Category \\
\midrule
EMA on landing prediction ($\alpha{=}0.5$)           & $-0.36$ & smoothing \\
$v_y$ median-warmup window after bounce              & $-1.54$ & smoothing \\
\texttt{opp\_y} idle pre-positioning                  & $-1.12$ & anticip. \\
Adaptive horizon                                     & neg.    & planning \\
Hybrid $\alpha$ raw velocity EMA                     & neg.    & smoothing \\
rMM offset feature vs $t_y$                          & neg.    & feature \\
Wall-coordinate fix                                  & neg.    & model \\
Ball VBGS $K{=}4$ retrain                            & neg.    & model \\
Aim-below offset (composition w/ anticipatory)       & $-0.61$ & composition \\
Paddle-centered defense                              & $-1.21$ & anticip. \\
$\ldots$                                             &         &           \\
\bottomrule
\end{tabular}
\end{table}

\paragraph{A pattern.}
Every intervention that traded responsiveness for noise reduction
lost. Every intervention that committed to a stale prediction lost.
The interventions that won were either fixes to wrong inputs
(bounce detection, wall clip) or replacements of stochastic
averaging by deterministic search (MPPI $\to$ tree).
% TODO: this might deserve a sentence framing it as a finding about
% high-noise low-horizon control settings generally, not just Pong.

% =========================================================================
\section{Reference Points, Not Competitors}\label{sec:comparison}
% =========================================================================

We position our numbers against the AXIOM Gameworld~10k results and
the Atari~100k sample-efficient RL leaderboard as \emph{reference
points}, not as direct competitors. Three differences make any
head-to-head claim ill-founded, and we state them up front before
showing any numbers:

\begin{itemize}
  \item \textbf{Input modality.} AXIOM, BBF, EfficientZero, and
        DreamerV3 consume raw pixels. Our agent consumes four RAM
        bytes (\texttt{RAM[49, 54, 51, 50]}) after $1/255$
        normalisation. This is a strictly easier perceptual front
        end: the agent is told the ball's $(x,y)$ and both
        paddles' $y$ exactly. The substantive computational work
        of the paper --- VBGS dynamics, EFE planning, rMM aim
        policy --- is downstream of perception, and our claims are
        about that downstream work. A pixel front end would be a
        natural extension and is listed as future work.
  \item \textbf{Game.} Gameworld Bounce~\cite{ref_axiom} is a
        Pong-\emph{like} task with custom physics and a custom
        opponent. \texttt{ALE/Pong-v5} is the canonical Atari
        title with the original ALE opponent. The two games are
        not the same task and should not be benchmarked against
        each other on a common scoreboard.
  \item \textbf{Sample regime.} AXIOM trains on 10\,000
        environment steps with no pretraining. Our VBGS is
        pretrained on 30\,000 offline frames of free play (no
        reward signal), then run online with rMM accumulation
        across rallies but no gradient steps anywhere. The
        appropriate axis for comparison is probably
        ``interactions-to-skill,'' not ``frames'' or
        ``parameters,'' and we do not know how to compute it
        cleanly across these very different pipelines.
\end{itemize}

With those caveats explicit, Table~\ref{tab:axiom} collects the
reference points.

\begin{table}[t]
\caption{Reference points. AXIOM and baselines on Bounce
(Gameworld 10k, 10\,000 env.\ steps, 10 seeds, pixel input). Our
agent on \texttt{ALE/Pong-v5} (RAM input, $N{=}100$ episodes,
30\,000-frame pretraining).
\textbf{These are different games on different inputs}; rows are
not directly comparable.}\label{tab:axiom}
\centering
\begin{tabular}{lrll}
\toprule
Agent                                 & Score                  & Input  & Domain \\
\midrule
AXIOM                                 & $+27 \pm 13$           & pixels & Gameworld Bounce \\
DreamerV3                             & $+14 \pm 16$           & pixels & Gameworld Bounce \\
BBF                                   & $-1 \pm 15$            & pixels & Gameworld Bounce \\
\midrule
\textbf{Ours} (closed-form AIF)       & $+1.94$                & RAM    & ALE/Pong-v5 \\
Rule-based baseline                   & $-14.4$                & RAM    & ALE/Pong-v5 \\
\bottomrule
\end{tabular}
\end{table}

\paragraph{What is, in our view, comparable.}
The shared substrate across our work and AXIOM ---
VBGS components, NIW conjugate updates, closed-form variational
inference, EFE planning --- means that the \emph{architectural}
comparisons carry directly: single-object vs object-centric,
pretraining on free play vs online structure growth and Bayesian
model reduction, deterministic tree search vs sampled rollouts.
The empirical lessons from this paper --- responsiveness beats
smoothing in low-horizon noisy control; bounce-aware velocity
tracking is decisive in projectile tasks; deterministic search
beats sampled averaging at short horizons --- should generalise to
the AXIOM regime, but we cannot claim they do without porting our
agent to Gameworld Bounce. We list that port as the most important
single piece of future work.

% =========================================================================
\section{Discussion and Outlook}\label{sec:discussion}
% =========================================================================

\paragraph{What the result is, framed honestly.}
The minimum closed-form AIF agent we could write, given 30\,000
frames of free play and no gradient steps, gets to a mean score of
$+1.94$ on \texttt{ALE/Pong-v5}. That is barely net-winning in
absolute terms, and we are not claiming otherwise. What is more
interesting than the absolute number is the decomposition: the
agent's skill is concentrated in two changes
(bounce-aware velocity, deterministic tree search) that are neither
architectural changes to the generative model nor algorithmic
changes to inference. Most of the ``Active Inference'' part of the
agent is doing its job before either change is made; both changes
are about feeding it cleaner inputs and querying it more decisively.

\paragraph{What did not work.}
A wider ball mixture ($K{=}4$) starved on per-component data; an
online ball VBGS regressed the score by $\approx -10$ via a
posterior-drift mechanism we do not yet understand; every smoothing
intervention we tried on the landing prediction lost; anticipatory
pre-positioning lost in every composition. We document these as
load-bearing constraints for follow-up work, not as cleanup we
deferred.

\paragraph{Future work.}
The single most important next step is to port the agent to
Gameworld Bounce~\cite{ref_axiom}, so that a head-to-head comparison
with AXIOM, BBF, and DreamerV3 becomes meaningful. Three further
directions: (i) replace pretraining with online structure growth
and Bayesian model reduction, closing the gap to AXIOM on
methodology; (ii) diagnose the online ball VBGS posterior-drift
regression that currently blocks fully online learning of dynamics;
(iii) add a pixel front end (an object detector or slot
attention) so that the RAM concession can be retired and the
comparison axis becomes purely about dynamics modelling and
planning, not perception.

\begin{credits}
\subsubsection{\ackname} [TODO --- funding sources, lab acknowledgments,
Anthropic Claude tooling acknowledgment if appropriate.]

\subsubsection{\discintname}
The authors have no competing interests to declare that are relevant
to the content of this article.
\end{credits}

% =========================================================================
% Bibliography
% =========================================================================
\begin{thebibliography}{8}

\bibitem{ref_axiom}
Heins, C., Van de Maele, T., Tschantz, A., Linander, H., Markovic, D.,
Salvatori, T., et al.: AXIOM: Learning to Play Games in Minutes with
Expanding Object-Centric Models. arXiv:2505.24784 (2025)

\bibitem{ref_atari100k}
Kaiser, {\L}., Babaeizadeh, M., Mi{\l}os, P., et al.: Model-Based
Reinforcement Learning for Atari. arXiv:1903.00374 (2019)

\bibitem{ref_efficientzero}
Ye, W., Liu, S., Kurutach, T., Abbeel, P., Gao, Y.: Mastering Atari
Games with Limited Data. NeurIPS 2021. arXiv:2111.00210

\bibitem{ref_bbf}
Schwarzer, M., Ceron, J.S.O., Courville, A., Bellemare, M.G.,
Agarwal, R., Castro, P.S.: Bigger, Better, Faster: Human-level
Atari with Human-level Efficiency. ICML 2023. arXiv:2305.19452

\bibitem{ref_dreamerv3}
Hafner, D., Pasukonis, J., Ba, J., Lillicrap, T.: Mastering Diverse
Domains through World Models. arXiv:2301.04104 (2023)

\bibitem{ref_fep}
Friston, K.: The Free-Energy Principle: A Unified Brain Theory?
Nature Reviews Neuroscience \textbf{11}, 127--138 (2010)

\bibitem{ref_aif_tutorial}
Da Costa, L., Parr, T., Sajid, N., Veselic, S., Neacsu, V., Friston,
K.: Active Inference on Discrete State-Spaces: A Synthesis. Journal
of Mathematical Psychology \textbf{99}, 102447 (2020)

\bibitem{ref_vbgs}
[TODO --- canonical VBGS reference; possibly Bishop 2006 ch.\ 10 for
variational Gaussian mixtures.]

\end{thebibliography}

\end{document}
